\documentclass[12pt]{article}
\usepackage[utf8]{inputenc}
\usepackage[T1]{fontenc}
\usepackage[margin=1in]{geometry}
\usepackage{setspace}
\usepackage{xcolor}
\usepackage[colorlinks=true, urlcolor=blue!60!black, linkcolor=blue!60!black]{hyperref}
\setlength{\parindent}{0em}
\setlength{\parskip}{0.6em}
\renewcommand{\baselinestretch}{1.0}
\begin{document}
\begin{center}
{\LARGE \textit{Bolling v. Sharpe}, 347 U.S. 497 (1954)\par}
\vspace{0.5em}
{\large Case Brief --- Two-Page Edition\par}
\end{center}

\section{Memory Jogger}
Because it would be ``unthinkable'' for the Constitution to forbid school segregation by the states but allow it by the federal government, the Fifth Amendment's Due Process Clause bars segregated schools in D.C., launching reverse incorporation (\href{https://supreme.justia.com/cases/federal/us/347/497/}{347 U.S. 497}).

\section{Detailed Case Facts}
In 1950, after the D.C. Board of Education refused to open the new Sousa Junior High as an integrated school, the Consolidated Parents Group brought eleven Black students, including Spottswood Bolling Jr., on opening day; despite empty classrooms, the principal turned them away because of their race. Because D.C. is a federal enclave, the Fourteenth Amendment's Equal Protection Clause did not apply, so the students relied on Fifth Amendment due process.

\section{Procedural History}
The district court dismissed under \textit{Carr v. Corning}. The Supreme Court granted certiorari before judgment to hear the case with \textit{Brown v. Board of Education} and decided both on May 17, 1954.

\section{Judicial Votes}
Unanimous, 9--0: Chief Justice Warren wrote for the Court, with no concurrences or dissents.

\section{Holding}
Racial segregation in D.C. public schools violates Fifth Amendment due process. Due process and equal protection are not mutually exclusive; racial classifications are suspect, and segregation, unrelated to any proper governmental objective, arbitrarily deprives Black children of liberty.

\section{Analysis of Opinions}
Warren conceded the Fifth Amendment has no equal protection clause but held that both guarantees stem from ``our American ideal of fairness.'' Relying on \textit{Korematsu} and \textit{Hi\-ra\-ba\-ya\-shi}, he treated race as ``constitutionally suspect,'' and he grounded the holding in a broad conception of liberty rather than in \textit{Brown}'s social science. Finally, given \textit{Brown}, a lesser duty on the federal government would be unthinkable. Later decisions, culminating in \textit{Adarand} (1995), turned this rationale into near-congruence of federal and state equal protection.

\section{Examples: Future Applications}
\textbf{Same-Side:} a federal benefit paid to widows but not widowers is invalid, as in \textit{Weinberger v. Wiesenfeld}. \textbf{Opposite-Side:} federal alienage limits on Medicare survive rational-basis review given Congress's immigration power (\textit{Mathews v. Diaz}). \textbf{Fence-Sitter:} excluding Puerto Rico residents from SSI pits anti-discrimination norms against the Territory Clause.

\section{Critique}
Originalists note that no one in 1791 read the Fifth Amendment to bar racial discrimination, and that the reading makes the Fourteenth Amendment's Equal Protection Clause redundant; some would rest the result on the Citizenship Clause. Progressives call the ``unthinkable'' rationale ipse dixit and argue that \textit{Adarand}'s congruence ignores Congress's remedial power under Section 5.

\section{Key Quotations}
\begin{quote}
``The `equal protection of the laws' is a more explicit safeguard of prohibited unfairness than `due process of law,' and therefore we do not imply that the two are always interchangeable phrases. But, as this Court has recognized, discrimination may be so unjustifiable as to be violative of due process.''
\end{quote}
\noindent --- Chief Justice Warren, Opinion of the Court, 347 U.S. at 499.

\end{document}
